\documentclass[10pt,a4paper]{extarticle}
\usepackage[a4paper,margin=10mm,top=9mm,bottom=11mm]{geometry}
% ===== Shared preamble: Class 12 Chemistry (project style, from physics-formula-sheet-source.tex) =====
\usepackage[T1]{fontenc}
\usepackage{lmodern}
\usepackage{amsmath,amssymb}
\usepackage{multicol}
\usepackage[most]{tcolorbox}
\usepackage{enumitem}
\usepackage{booktabs,tabularx,array}
\usepackage{fancyhdr}
\usepackage{xcolor}
\usepackage{microtype}
\usepackage[version=4]{mhchem}
\usepackage{tikz}
\usetikzlibrary{arrows.meta,calc}
\usepackage{pifont}
\usepackage{needspace}
\DeclareUnicodeCharacter{21CC}{\ensuremath{\rightleftharpoons}}
\DeclareUnicodeCharacter{2713}{\ding{51}}
\DeclareUnicodeCharacter{03B1}{\ensuremath{\alpha}}
\DeclareUnicodeCharacter{03B2}{\ensuremath{\beta}}
\DeclareUnicodeCharacter{2192}{\ensuremath{\to}}
\DeclareUnicodeCharacter{2610}{\ensuremath{\square}}
\DeclareUnicodeCharacter{21D2}{\ensuremath{\Rightarrow}}
\DeclareUnicodeCharacter{2194}{\ensuremath{\leftrightarrow}}
\DeclareUnicodeCharacter{2605}{\ding{72}}

% ---------- Colours: one per unit ----------
\definecolor{navy}{HTML}{1B2A4A}
\definecolor{uI}{HTML}{1F5FAD}    % 1 Solutions
\definecolor{uII}{HTML}{0E8A7E}   % 2 Electrochemistry
\definecolor{uIII}{HTML}{C2571A}  % 3 Chemical Kinetics
\definecolor{uIV}{HTML}{6A3FA0}   % 4 d- and f-Block
\definecolor{uV}{HTML}{B0204F}    % 5 Coordination Compounds
\definecolor{uVI}{HTML}{2E8B3A}   % 6 Haloalkanes & Haloarenes
\definecolor{uVII}{HTML}{8A6D00}  % 7 Alcohols, Phenols, Ethers
\definecolor{uVIII}{HTML}{2B6F8F} % 8 Aldehydes, Ketones, Acids
\definecolor{uIX}{HTML}{9A4A1E}   % 9 Amines
\definecolor{uX}{HTML}{5A5A5A}    % 10 Biomolecules
\definecolor{tipY}{HTML}{FFF4B8}
\definecolor{tipYframe}{HTML}{E0B400}
\definecolor{sqp}{HTML}{1B2A4A}
\colorlet{uc}{uI}

% ---------- Tag for questions in the official 2026-27 sample paper ----------
\newcommand{\tagbox}[2]{{\setlength{\fboxsep}{0.8pt}\colorbox{#1}{\textcolor{white}{\sffamily\bfseries\tiny\,#2\,}}}}
\newcommand{\SQP}{\tagbox{sqp}{SQP 26--27}}

% ---------- Unit banner & chapter bar ----------
\newcommand{\unitbanner}[3]{% colour, title, marks
  \colorlet{uc}{#1}%
  \par\vspace{2pt}\noindent
  \begin{tcolorbox}[enhanced,colback=#1,colframe=#1,arc=1.5pt,boxsep=1pt,left=4pt,right=4pt,top=1.5pt,bottom=1.5pt,before skip=2pt,after skip=2pt]
  \sffamily\bfseries\color{white}\footnotesize #2\hfill\scriptsize #3
  \end{tcolorbox}}
\newcommand{\chapbar}[1]{%
  \par\needspace{7\baselineskip}\noindent
  \begin{tcolorbox}[enhanced,colback=uc!18,colframe=uc,boxrule=0pt,leftrule=3pt,arc=0pt,boxsep=0.5pt,left=3pt,right=3pt,top=1pt,bottom=1pt,before skip=3pt,after skip=1.5pt]
  \sffamily\bfseries\color{uc!80!black}\footnotesize #1
  \end{tcolorbox}}
\newcommand{\sub}[1]{\par\needspace{4\baselineskip}\addvspace{2pt}\noindent{\sffamily\bfseries\scriptsize\color{uc}\strut\MakeUppercase{#1}}\par\vspace{0.5pt}}

% ---------- Boxes ----------
\newtcolorbox{kbox}{enhanced,colback=uc!5,colframe=uc,boxrule=0.6pt,arc=1.5pt,left=2pt,right=2pt,top=1pt,bottom=1pt,boxsep=1pt,before skip=2pt,after skip=2pt}
\newtcolorbox{tip}{enhanced,colback=tipY,colframe=tipYframe,boxrule=0.5pt,arc=1.5pt,left=2pt,right=2pt,top=0.5pt,bottom=0.5pt,boxsep=1pt,before skip=2pt,after skip=2pt,
  fontupper=\footnotesize}
\newtcolorbox{numbox}{enhanced,colback=white,colframe=uc!70,boxrule=0.5pt,arc=1.5pt,left=2pt,right=2pt,top=0.5pt,bottom=0.5pt,boxsep=1pt,before skip=2pt,after skip=2pt,
  borderline={0.4pt}{1.2pt}{uc!35,dashed}}
\newcommand{\Tip}{{\sffamily\bfseries\scriptsize TIP\;}}
\newcommand{\Num}{{\sffamily\bfseries\scriptsize\color{uc}STANDARD NUMERICALS\;}}

% ---------- Lists ----------
\setlist{nosep,leftmargin=8pt,labelsep=2pt}
\setlist[itemize]{label={\textcolor{uc}{\tiny$\blacktriangleright$}}}
\newenvironment{fl}{\begin{itemize}}{\end{itemize}}

% ---------- Mini graph helper (tikz) ----------
% \minigraph{width}{height}{x label}{y label}{drawing commands}
\newcommand{\minigraph}[5]{%
\begin{tikzpicture}[x=1cm,y=1cm,line cap=round,font=\tiny]
\draw[-{Stealth[length=3pt]},black!70] (0,0)--(#1,0) node[below left,inner sep=1pt]{#3};
\draw[-{Stealth[length=3pt]},black!70] (0,0)--(0,#2) node[above right,inner sep=1pt]{#4};
#5
\end{tikzpicture}}

\setlength{\columnsep}{9pt}
\setlength{\columnseprule}{0.3pt}
\def\columnseprulecolor{\color{black!25}}
\setlength{\parindent}{0pt}
\setlength{\abovedisplayskip}{1pt}\setlength{\belowdisplayskip}{1pt}
\setlength{\abovedisplayshortskip}{0pt}\setlength{\belowdisplayshortskip}{0pt}
\renewcommand{\arraystretch}{1.08}
\newcolumntype{Y}{>{\raggedright\arraybackslash}X}
\pagestyle{fancy}\fancyhf{}
\renewcommand{\headrulewidth}{0pt}

\fancyfoot[L]{\tiny\sffamily\color{black!55} Class 12 Chemistry · \sheetname · CBSE Board 2026--27 (paper 27 Feb 2027)}
\fancyfoot[R]{\tiny\sffamily\color{black!55} Page \thepage}

% ---------- Notes-specific macros ----------
\newcommand{\st}[1]{\textcolor{tipYframe!90!black}{\ifcase#1\or\ding{72}\or\ding{72}\ding{72}\or\ding{72}\ding{72}\ding{72}\fi}}
\newcommand{\ty}[1]{\tagbox{uc!75!black}{#1}}
\newcommand{\mk}[1]{{\scriptsize\color{black!55}[#1]}}
\newenvironment{rqlist}{\begin{enumerate}[leftmargin=12pt,label={\sffamily\bfseries\scriptsize\color{uc}\arabic*.},itemsep=1pt]}{\end{enumerate}}
\newcommand{\ans}{\par\hspace*{2pt}{\color{uc}$\hookrightarrow$}\;}
\newtcolorbox{skel}[1]{enhanced,colback=white,colframe=uc,boxrule=0.5pt,arc=1.5pt,left=3pt,right=3pt,top=1pt,bottom=1pt,boxsep=1pt,before skip=2pt,after skip=3pt,
  title={\sffamily\bfseries\scriptsize #1},colbacktitle=uc!85,coltitle=white,toptitle=0.5pt,bottomtitle=0.5pt}
\newcommand{\chapterhead}[4]{% colour, title, marks, subtitle
  \newpage\colorlet{uc}{#1}%
  \begin{tcolorbox}[enhanced,colback=#1,colframe=#1,arc=2pt,left=6pt,right=6pt,top=3pt,bottom=3pt,after skip=3pt]
  \color{white}\sffamily{\large\bfseries #2}\hfill{\small\bfseries #3}\\[-1pt]{\scriptsize #4}
  \end{tcolorbox}}

% orbital boxes: 0 empty, 1 up, 2 pair, 3 ligand electron pair (shaded)
\newcommand{\ob}[1]{\begin{tikzpicture}[x=2.9mm,y=2.9mm,baseline=0.3mm,line width=0.35pt]
\foreach \s [count=\i from 0] in {#1}{%
 \ifnum\s=3 \fill[uc!28] (\i,0) rectangle ++(1,1);\fi
 \draw[black!65] (\i,0) rectangle ++(1,1);
 \ifnum\s>0 \draw[-{Stealth[length=1.6pt]}] (\i+0.33,0.12)--(\i+0.33,0.88);\fi
 \ifnum\s>1 \draw[{Stealth[length=1.6pt]}-] (\i+0.67,0.12)--(\i+0.67,0.88);\fi}
\end{tikzpicture}}
\def\sheetname{Expert Q\&A}
\raggedbottom
\begin{document}
\small

% ---------- Q&A sheet macros (used after \begin{document}) ----------
\newcounter{qn}\newcounter{an}
\newcommand{\Q}{\par\addvspace{4pt}\refstepcounter{qn}\noindent\hangindent=21pt\hangafter=1\makebox[21pt][l]{\sffamily\bfseries\footnotesize\color{uc}Q\theqn.}}
\newcommand{\An}{\par\addvspace{4pt}\refstepcounter{an}\noindent\hangindent=21pt\hangafter=1\makebox[21pt][l]{\sffamily\bfseries\footnotesize\color{uc}A\thean.}}
\newcommand{\nl}{\par\noindent\hangindent=21pt\hangafter=0}
\newcommand{\secbar}[2]{\par\addvspace{4pt}\needspace{5\baselineskip}\noindent
  \begin{tcolorbox}[enhanced,colback=uc!18,colframe=uc,boxrule=0pt,leftrule=3pt,arc=0pt,boxsep=0.5pt,left=3pt,right=3pt,top=1pt,bottom=1pt,before skip=3pt,after skip=2pt]
  \sffamily\bfseries\color{uc!80!black}\footnotesize #1\hfill\scriptsize #2
  \end{tcolorbox}}
\newcommand{\qatitle}[5]{% number, title, chapters, marks, time
  \gdef\sheetname{Expert Q\&A Sheet #1 of 6: #2}
  \begin{tcolorbox}[enhanced,colback=navy,colframe=navy,arc=2pt,left=8pt,right=8pt,top=5pt,bottom=5pt]
  \color{white}\sffamily
  {\LARGE\bfseries Class 12 Chemistry — Expert Q\&A Sheet #1 of 6}\\[3pt]
  {\large #2} \hfill {\small #3}\\[2pt]
  {\footnotesize CBSE Board 2026--27 pattern · Rationalised NCERT · Paper: \textbf{Sat 27 Feb 2027} · \textbf{#4 marks · #5 min} · Questions first, full-mark answers with step marks at the end}
  \end{tcolorbox}
  \begin{tcolorbox}[enhanced,colback=tipY,colframe=tipYframe,boxrule=0.5pt,arc=1.5pt,left=4pt,right=4pt,top=1.5pt,bottom=1.5pt,before skip=3pt,after skip=3pt]
  \footnotesize\textbf{How to use.} Sit it like a test: set a timer, write on paper, no notes. Then mark yourself against the answer section, step by step. \textbf{Honest label:} every question is \textbf{PYQ-style} — the question types and data patterns that recur in CBSE papers of 2022--2026 and in NCERT — written in my own words, not copied from a specific year's paper. Step-mark splits follow the usual CBSE marking-scheme pattern and are estimates.
  \end{tcolorbox}}
\newcommand{\answerstart}{\newpage
  \begin{tcolorbox}[enhanced,colback=navy,colframe=navy,arc=2pt,left=6pt,right=6pt,top=3pt,bottom=3pt,after skip=3pt]
  \color{white}\sffamily{\large\bfseries Answers \& marking scheme}\hfill{\small written the way a full-marks script reads · {\scriptsize[\,]} = step marks}
  \end{tcolorbox}}
\newcommand{\scorebox}[1]{%
  \begin{tcolorbox}[enhanced,colback=white,colframe=navy,boxrule=0.8pt,arc=2pt,title={\sffamily\bfseries\footnotesize My score: \underline{\hspace{10mm}} / #1 \quad · \quad Time taken: \underline{\hspace{10mm}} min \quad · \quad Questions to redo / errors to log},colbacktitle=navy,coltitle=white,height fill,before skip=4pt]
  \end{tcolorbox}}

\colorlet{uc}{uVI}
\qatitle{3}{Organic Chemistry A}{Ch 6 Haloalkanes and Haloarenes · Ch 7 Alcohols, Phenols and Ethers}{42}{75}
\begin{multicols}{2}\raggedcolumns
\secbar{SECTION A · MCQ and assertion--reason}{8 × 1 = 8}
\Q Which reacts fastest by the S$_\text{N}$1 mechanism?\nl (a) \ce{CH3Br} \quad (b) \ce{CH3CH2Br} \quad (c) \ce{(CH3)2CHBr} \quad (d) \ce{(CH3)3CBr} \mk{1}
\Q The main product of bromoethane with AgCN is\nl (a) propanenitrile \quad (b) ethyl isocyanide \quad (c) ethyl nitrite \quad (d) nitroethane \mk{1}
\Q Which molecule is chiral?\nl (a) 2-bromopropane \quad (b) 2-bromobutane \quad (c) 1-bromobutane \quad (d) 2-bromo-2-methylpropane \mk{1}
\Q Which alcohol gives turbidity \emph{immediately} with Lucas reagent?\nl (a) butan-1-ol \quad (b) butan-2-ol \quad (c) 2-methylpropan-2-ol \quad (d) 2-methylpropan-1-ol \mk{1}
\Q Phenol reacts with bromine water to give\nl (a) 2-bromophenol \quad (b) 4-bromophenol \quad (c) 2,4-dibromophenol \quad (d) 2,4,6-tribromophenol \mk{1}
\Q The strongest acid among the following is\nl (a) phenol \quad (b) 4-methylphenol \quad (c) 4-nitrophenol \quad (d) ethanol \mk{1}
\nl\textit{For Q7--Q8 choose: (a) A and R true, R explains A; (b) both true, R does not explain A; (c) A true, R false; (d) A false, R true.}
\Q \textbf{A:} Ethanol boils at a higher temperature than methoxymethane. \textbf{R:} Ethanol molecules are associated through intermolecular hydrogen bonds. \mk{1}
\Q \textbf{A:} Chlorobenzene is more reactive than chloromethane towards nucleophilic substitution. \textbf{R:} The C--Cl bond in chlorobenzene has partial double-bond character. \mk{1}

\secbar{SECTION B · Very short answer}{4 × 2 = 8}
\Q Write one equation each for (i) the Finkelstein reaction, (ii) the Wurtz--Fittig reaction. \mk{2}
\Q In each pair, which compound reacts faster, and why?\nl (i) S$_\text{N}$2: 1-bromobutane or 2-bromobutane\nl (ii) S$_\text{N}$1: \ce{C6H5CH2Cl} or \ce{C6H5CH(Cl)C6H5} \mk{2}
\Q Write the equations for (i) Kolbe's reaction, (ii) the Reimer--Tiemann reaction. \mk{2}
\Q Give reasons: (i) o-nitrophenol is steam-volatile but p-nitrophenol is not; (ii) Grignard reagents must be prepared under anhydrous conditions. \mk{2}

\secbar{SECTION C · Short answer}{4 × 3 = 12}
\Q Write the mechanism of the acid-catalysed hydration of ethene to ethanol. \mk{3}
\Q How will you carry out these conversions?\nl (i) but-1-ene to but-2-ene \quad (ii) benzene to biphenyl \quad (iii) cumene to phenol \mk{3}
\Q (a) Give two reasons why haloarenes are much less reactive than haloalkanes towards nucleophilic substitution. \mk{2}\nl (b) Name the major alkene formed when 2-bromopentane is heated with alcoholic KOH, and the rule that predicts it. \mk{1}
\Q (a) Write the products of (i) \ce{(CH3)3C-O-CH3 + HI}, (ii) \ce{C6H5-O-CH3 + HI}. \mk{2}\nl (b) Why can't tert-butyl methyl ether be made from \ce{(CH3)3C-Br} and \ce{CH3ONa}? \mk{1}

\secbar{SECTION D · Case-based}{1 × 4 = 4}
\Q \textit{Nucleophilic substitution of haloalkanes follows two main pathways. In S$_\text{N}$2 the nucleophile attacks the carbon from the side opposite the leaving group in a single step. In S$_\text{N}$1 the C--X bond first breaks to give a carbocation, which is then attacked by the nucleophile. The stereochemical result, studied with optically active haloalkanes, is different in the two cases.}\\[2pt]
(i) What happens to the configuration when ($-$)-2-bromooctane reacts with NaOH by the S$_\text{N}$2 path? \mk{1}\\
(ii) Why does an optically active haloalkane give a racemic product by the S$_\text{N}$1 path? \mk{1}\\
(iii) Arrange 1-bromopentane, 2-bromopentane and 2-bromo-2-methylbutane in increasing order of S$_\text{N}$2 reactivity, with a reason. \mk{2}\\
\textbf{OR} (iii) What is a chiral molecule? Which of 1-chlorobutane and 2-chlorobutane is chiral? What is a racemic mixture? \mk{2}

\secbar{SECTION E · Long answer}{2 × 5 = 10}
\Q Compound A (\ce{C3H8O}) gives turbidity with Lucas reagent in about five minutes. When vapours of A are passed over copper at 573 K, B (\ce{C3H6O}) is formed. Heating A with conc.\ \ce{H2SO4} at 443 K gives C (\ce{C3H6}). C adds HBr to give D, and D gives back A with aqueous KOH. D reacts with sodium in dry ether to give E (\ce{C6H14}). Identify A to E and write the equations. \mk{5}
\Q (a) Arrange ethanol, phenol, 4-methylphenol and 4-nitrophenol in increasing order of acid strength and explain the order. \mk{2}\\
(b) Write equations for what happens when: (i) phenol is heated with zinc dust; (ii) propan-1-ol is oxidised with PCC; (iii) anisole is treated with \ce{CH3Cl} and anhydrous \ce{AlCl3}. \mk{3}
\end{multicols}

\answerstart
\begin{multicols}{2}\raggedcolumns
\secbar{SECTION A}{1 mark each}
\An \textbf{(d) \ce{(CH3)3CBr}.} It forms the most stable (3$^\circ$) carbocation.
\An \textbf{(b) ethyl isocyanide,} \ce{CH3CH2NC}. AgCN is mainly covalent, so nitrogen is the donor atom.
\An \textbf{(b) 2-bromobutane.} C-2 carries four different groups: H, Br, \ce{CH3}, \ce{C2H5}.
\An \textbf{(c) 2-methylpropan-2-ol} (a 3$^\circ$ alcohol).
\An \textbf{(d) 2,4,6-tribromophenol} (white precipitate).
\An \textbf{(c) 4-nitrophenol.} \ce{-NO2} withdraws electrons and stabilises the phenoxide ion.
\An \textbf{(a).}
\An \textbf{(d).} Chlorobenzene is \emph{less} reactive; the reason given is the correct explanation of that.

\secbar{SECTION B}{2 marks each}
\An (i) \ce{CH3CH2Br + NaI ->[dry acetone] CH3CH2I + NaBr} \mk{1}\nl
(ii) \ce{C6H5Cl + CH3Cl + 2Na ->[dry ether] C6H5CH3 + 2NaCl} \mk{1}
\An (i) \textbf{1-Bromobutane}: it is a primary halide, so there is less steric hindrance to backside attack by the nucleophile. \mk{1}\nl
(ii) \textbf{\ce{C6H5CH(Cl)C6H5}}: its carbocation is stabilised by resonance with two phenyl rings, so it forms more easily. \mk{1}
\An (i) Sodium phenoxide $\xrightarrow[\text{(ii) H}^+]{\text{(i) CO}_2,\ 400\text{ K},\ 4\text{--}7\text{ atm}}$ 2-hydroxybenzoic acid (salicylic acid) \mk{1}\nl
(ii) Phenol $\xrightarrow[\text{(ii) H}^+]{\text{(i) CHCl}_3\text{ + aq. NaOH},\ 340\text{ K}}$ 2-hydroxybenzaldehyde (salicylaldehyde) \mk{1}\nl
{\scriptsize In the exam, draw the benzene rings with --OH and the new group at the ortho position.}
\An (i) o-Nitrophenol has \emph{intra}molecular hydrogen bonding, so its molecules are not associated and it is volatile; p-nitrophenol has \emph{inter}molecular hydrogen bonding, so its molecules are associated. \mk{1}\nl
(ii) Grignard reagents react with water (any proton source) to give alkanes: \ce{RMgX + H2O -> RH + Mg(OH)X}. \mk{1}

\secbar{SECTION C}{3 marks each}
\An \textbf{Step 1} — protonation of the alkene by \ce{H3O+} to form a carbocation:\nl
\ce{CH2=CH2 + H3O+ -> CH3-CH2+ + H2O} \mk{1}\nl
\textbf{Step 2} — nucleophilic attack of water on the carbocation:\nl
\ce{CH3-CH2+ + H2O -> CH3-CH2-OH2+} \mk{1}\nl
\textbf{Step 3} — deprotonation to give the alcohol:\nl
\ce{CH3-CH2-OH2+ + H2O -> CH3CH2OH + H3O+} \mk{1}
\An (i) \ce{CH3CH2CH=CH2 ->[HBr] CH3CH2CH(Br)CH3 ->[alc. KOH][\Delta] CH3CH=CHCH3} \mk{1}\nl
(ii) \ce{C6H6 ->[Cl2][anhyd. FeCl3] C6H5Cl};\; \ce{2C6H5Cl + 2Na ->[dry ether] C6H5-C6H5 + 2NaCl} \mk{1}\nl
(iii) Cumene \ce{->[O2]} cumene hydroperoxide \ce{->[H+][H2O]} phenol + propanone \mk{1}
\An (a) (i) \textbf{Resonance}: a lone pair of the halogen is delocalised into the ring, giving the C--X bond partial double-bond character, so it is harder to break. \mk{1}\nl
(ii) \textbf{Hybridisation}: the carbon bonded to halogen is sp$^2$ (more electronegative than sp$^3$), so the C--X bond is shorter and stronger. (Also acceptable: the phenyl cation is not resonance-stabilised; the electron-rich ring repels the nucleophile.) \mk{1}\nl
(b) \textbf{Pent-2-ene} is the major product, by \textbf{Zaitsev's rule} (the alkene with more alkyl groups on the doubly bonded carbons is preferred). \mk{1}
\An (a) (i) \ce{(CH3)3C-O-CH3 + HI -> (CH3)3C-I + CH3OH} (S$_\text{N}$1 through the stable 3$^\circ$ carbocation) \mk{1}\nl
(ii) \ce{C6H5-O-CH3 + HI -> C6H5OH + CH3I} (the O--aryl bond has partial double-bond character and is not broken) \mk{1}\nl
(b) Methoxide is a strong base; with a 3$^\circ$ halide, elimination wins over substitution and 2-methylpropene is formed. The ether is made from \ce{(CH3)3C-ONa + CH3Br} instead. \mk{1}

\secbar{SECTION D}{4 marks}
\An (i) The configuration is \textbf{inverted}: the product is (+)-octan-2-ol, with the --OH on the side opposite to where Br was. \mk{1}\nl
(ii) The carbocation intermediate is planar (sp$^2$); the nucleophile attacks from either face with equal ease, giving both enantiomers in equal amounts. \mk{1}\nl
(iii) 2-bromo-2-methylbutane (3$^\circ$) $<$ 2-bromopentane (2$^\circ$) $<$ 1-bromopentane (1$^\circ$). \mk{1} More alkyl groups on the carbon bearing Br mean more steric hindrance to the approaching nucleophile. \mk{1}\nl
\textbf{OR} (iii) A chiral molecule is one that is not superimposable on its mirror image. \mk{½} \textbf{2-Chlorobutane} is chiral (C-2 has four different groups). \mk{½} A racemic mixture contains equal amounts of the two enantiomers and has zero optical rotation. \mk{1}

\secbar{SECTION E}{5 marks each}
\An Turbidity in about five minutes ⇒ A is a secondary alcohol.\nl
\textbf{A = propan-2-ol}, \ce{CH3CH(OH)CH3} \quad \textbf{B = propanone}, \ce{CH3COCH3}\nl
\textbf{C = propene}, \ce{CH3CH=CH2} \quad \textbf{D = 2-bromopropane}, \ce{CH3CH(Br)CH3}\nl
\textbf{E = 2,3-dimethylbutane}, \ce{(CH3)2CH-CH(CH3)2} \mk{½ each = 2½}\nl
\ce{CH3CH(OH)CH3 ->[Cu][573 K] CH3COCH3 + H2} \mk{½}\nl
\ce{CH3CH(OH)CH3 ->[conc. H2SO4][443 K] CH3CH=CH2 + H2O} \mk{½}\nl
\ce{CH3CH=CH2 + HBr -> CH3CH(Br)CH3} (Markovnikov) \mk{½}\nl
\ce{CH3CH(Br)CH3 + KOH(aq) -> CH3CH(OH)CH3 + KBr} \mk{½}\nl
\ce{2CH3CH(Br)CH3 + 2Na ->[dry ether] (CH3)2CH-CH(CH3)2 + 2NaBr} \mk{½}
\An (a) ethanol $<$ 4-methylphenol $<$ phenol $<$ 4-nitrophenol \mk{1}\nl
Phenols are more acidic than ethanol because the phenoxide ion is stabilised by resonance; ethoxide is not. \ce{-NO2} ($-I$, $-R$) spreads the negative charge further and raises acidity; \ce{-CH3} (+I) intensifies the charge and lowers it. \mk{1}\nl
(b) (i) \ce{C6H5OH + Zn ->[\Delta] C6H6 + ZnO} \mk{1}\nl
(ii) \ce{CH3CH2CH2OH ->[PCC] CH3CH2CHO} (propanal; oxidation stops at the aldehyde) \mk{1}\nl
(iii) \ce{C6H5OCH3 + CH3Cl ->[anhyd. AlCl3]} 4-methoxytoluene (major) + 2-methoxytoluene (minor) + HCl \mk{1}
\end{multicols}
\begin{tcolorbox}[enhanced,colback=tipY,colframe=tipYframe,boxrule=0.5pt,arc=1.5pt,left=4pt,right=4pt,top=2pt,bottom=2pt,before skip=4pt]
\footnotesize\textbf{Where a 70/70 script wins marks here.} (1) Every arrow carries its reagent \emph{and} condition (dry ether, dry acetone, anhyd.\ \ce{AlCl3}, 443 K). (2) In ``which is faster'' questions the answer names the compound first, then gives one structural reason. (3) Mechanisms are written as numbered steps with the intermediate shown. (4) In ``identify A--E'' questions, each letter gets a name \emph{and} a structure, and each change gets an equation. (5) For ring compounds, draw the structures in the exam; a name alone can lose half a mark.
\end{tcolorbox}
\scorebox{42}

\end{document}
